\section{Related Works}
\label{sec:related_works}

Optimization-based control has become a standard approach to dynamic legged
locomotion. It formulates motion generation as an Optimal Control Problem in
which the robot dynamics, contact forces, physical constraints, and task
objectives are considered explicitly. Solving the complete problem online is
difficult because legged robots have many degrees of freedom and interact with
the environment through intermittent and unilateral contacts
\cite{legged_optimization}. Existing controllers therefore differ in the
dynamic model they employ, the way they represent contacts, and the division of
tasks between predictive planning and low-level control.\\
\
Reduced-order models make online optimization tractable by retaining the
variables that dominate balance and locomotion while neglecting part of the
articulated dynamics. The convex MPC developed for the MIT Cheetah 3 represents
the trunk as a single rigid body and optimizes the ground reaction forces over
a finite horizon \cite{cheetah_mpc}. The whole-body nonlinear MPC of Neunert
et al. instead includes the complete articulated dynamics and explicit contact
dynamics, optimizing robot motion,
contact locations, sequences, and timings within a single problem
\cite{whole_body_nmpc}. This richer representation
reduces the separation between planning and control but results in a more
demanding optimization. Other work has addressed the numerical treatment of
three-dimensional orientation \cite{representation_free_mpc} and adapted solver
accuracy to the time constraints of whole-body locomotion control
\cite{fast_whole_body_mpc}.\\
\
Wheeled-legged robots introduce rolling contacts and additional actuated
degrees of freedom. Early experiments on ANYmal integrated wheel motion and
non-holonomic rolling constraints into a whole-body planning and control
framework \cite{keep_rollin}. Online trajectory optimization was later used to
produce hybrid walking-driving motions on rough terrain
\cite{rolling_deep}. A subsequent whole-body MPC formulation jointly optimize
torso motion, wheel motion, and contact forces using a kinodynamic model that
combines Single Rigid Body Dynamics with whole-body kinematics. Unlike the
full-dynamics formulation discussed above, this model does not retain the
complete articulated dynamics. A separate gait-sequence generator uses
kinematic leg utilities to adapt contact timings online
\cite{wheeled_whole_body_mpc}. These approaches coordinate wheel motion with
leg configuration, body dynamics, and contact forces rather than treating the
wheels as an independent drivetrain.\\
\
The wheeled-legged studies discussed above focus on terrain traversal and the
transition between rolling and stepping. At higher speeds, vehicle-like
effects become increasingly relevant. In particular, lateral load transfer
during cornering can unload the inner wheels and bring the robot closer to
rollover \cite{racing}. A recent controller for the Unitree
Go2-W combined an MPC planner for active roll regulation with a learned
low-level policy. The legs were used as an active suspension system to lean the
body into turns and reduce lateral load transfer during physical racing
experiments \cite{racing}.\\
\
Reinforcement Learning provides a different approach to the uncertainties and
contact interactions involved in locomotion. Policies trained in simulation
have been transferred to physical quadrupeds for agile movements and locomotion
over challenging terrain \cite{agile_motor_skills,challenging_terrain}.
Perceptive policies have also demonstrated locomotion in natural environments
using proprioceptive and exteroceptive information
\cite{perceptive_rl}. These methods can learn behaviours that are difficult to
describe through an analytical model, but their performance depends on the
training distribution, reward formulation, observation design, and simulation
accuracy.\\
\
The computational structure of RL has encouraged massively parallel
simulation. Thousands of robots can be simulated simultaneously on a GPU,
producing large batches of experience for policy optimization. This approach
has substantially reduced the time required to train locomotion policies and
supports curriculum learning across many parallel environments
\cite{parallel_rl}. A purely model-free policy, however, does not explicitly
encode the physical constraints or the structure available to an
optimization-based controller. This limitation has motivated methods that
combine learning and model-based control.\\
\
Some hybrid architectures combine a model-based motion planner with an RL
policy trained to execute the planned motion. In \cite{deep_tracking},
trajectory optimization generates footholds and base motion and, during training,
rewards encourage base-motion tracking and accurate foot placement, while the
actor observes selected planning outputs rather than the full optimized base
trajectory. The planner remains active during deployment.
Lee et al. instead use Linear Inverted Pendulum dynamics
to generate desired foot placements from velocity commands. Their RL policy
tracks those placements without being required to follow the complete motion
predicted by the reduced model \cite{footstep_rl}. In both cases, the model
guides the learned behaviour, while the policy remains responsible for
producing the low-level control action.\\
\
Hierarchical methods  can also assign high-level decisions to RL and motion generation
to MPC, reversing the roles of the model-based and learned components in the
preceding examples. In
\cite{non_gaited}, a high-level RL policy selects acyclic contact patterns and
navigation commands for a low-level MPC. This removes the combinatorial search
over contact sequences from the continuous optimization and allows the same
architecture to address legged and hybrid wheeled-legged locomotion.\\
\
Other approaches use learning to adapt the prediction model or individual
components of the nominal control architecture. Learned corrections have been introduced
to reduce the
discrepancy between reduced-order predictions and the full-order dynamics of a
quadruped \cite{model_gap}. Modular residual methods apply learned corrections
to individual parts of a model-based controller, including its dynamics model
and footstep planner \cite{modular_residual}. In the adaptive architecture
presented in \cite{adaptive_mpc}, RL modifies the predictive dynamics, the
swing-leg controller, and the gait frequency. These adaptations affect both
the MPC problem and the generation of swing-leg commands, rather than adding
a residual torque to the nominal actuator command.\\
\
Residual corrections can therefore act inside the model-based architecture,
as in the model and planner adaptations discussed above. Action-level residual
policy learning instead applies the correction to the nominal control output. A
conventional
controller supplies a nominal action, while RL learns only the correction
needed to improve the resulting behaviour. This approach has been demonstrated
in contact-rich manipulation tasks, where a feedback controller captured the
known structure and a residual policy compensated for effects that were
difficult to model \cite{residual_robot_control}. For locomotion, corrections
to joint setpoints and corrections expressed directly as torques provide
different ways of modifying the nominal actuator command
\cite{residual_mpc}.\\
\
The work most closely related to this thesis is Residual MPC, which executes an
MPC controller and a residual policy concurrently and combines their outputs at
the torque-control level \cite{residual_mpc}. The MPC provides a structured
control prior, while the policy learns corrections for modelling errors and
behaviours that are difficult to encode in the optimization. The paper compares
different residual action spaces and torque-blending strategies, showing that
the action representation and the way the outputs are combined affect both
learning and final locomotion behaviour \cite{residual_mpc}. Unlike approaches based on
offline demonstrations, MPC
remains active during training and deployment.\\
\
This integration creates a substantial computational requirement. An MPC
problem must be solved for every simulated robot at each control update.
Optimal-control solvers for legged locomotion have traditionally been designed
for sequential execution on CPUs, whereas modern RL environments simulate
thousands of agents in parallel on GPUs. The resulting separation between CPU
optimization and GPU training limits the number of MPC instances that can be
kept inside the learning loop \cite{residual_mpc,mpx}.\\
\
CusADi addressed this problem by translating symbolic expressions into
GPU-compatible operations and evaluating large batches of approximate
optimal-control problems in parallel \cite{cusadi}. ReLU-QP addresses the
solution of quadratic programs by expressing ADMM iterations as a weight-tied
network with ReLU activations, allowing execution through standard
machine-learning tools on GPUs \cite{relu_qp}. MPX uses a primal-dual iLQR formulation
implemented in JAX. Associative scans enable temporal and state-space
parallelization, while batching supports many independent MPC instances on
the same device \cite{mpx}. This enables MPC to remain directly inside a
massively parallel RL training loop.\\
\
These studies motivate the combination of a model-based locomotion controller
with learned corrections, but their models and control hierarchies cannot be
transferred unchanged to TITA. The robot shares the rolling-contact constraints
of wheeled quadrupeds, while its two-wheel support makes body balance and the
coordination of wheel traction with leg motion central to the control problem.
This thesis adopts the action-level residual approach: a DFCIP-based MPC and a
constrained Whole-Body Controller generate the nominal command, while the
policy modifies torque-level references to produce an additional torque
contribution. The MPC and WBC remain active during training. Compared with the
kinodynamic whole-body MPC used in Residual MPC \cite{residual_mpc}, the TITA
architecture separates reduced-order prediction from whole-body torque
computation. The architecture chapter describes this decomposition and how
the learned correction is combined with the nominal command.
